\section{Gedächtnis Teil 1}
\slide{Begrifflichkeiten}{
    \begin{itemize}
        \item \b{Rehearsal:} Der Prozess der aktiven Wiederholung von Informationen (z. B. durch eine "innere Stimme"), um sie im Kurzzeitgedächtnis zu halten und das Vergessen zu verzögern.
        \item \b{Chunk:} Eine sinnvolle Informationseinheit, die durch das Zusammenfassen einzelner Elemente (z. B. Ziffern zu einem Datum) gebildet wird, um die Gedächtnisspanne zu erweitern.
        \item \b{Retrieval Structure:} Eine mentale Struktur (z. B. ein Schema), die hilft, Informationen beim Abruf wieder zu entpacken ("Ent-Chunking") und zu organisieren.
        \item \b{Millersche Zahl:} Beschreibt die begrenzte Kapazität des Kurzzeitgedächtnisses, die bei etwa $7 \pm 2$ Einheiten (Chunks) liegt.
    \end{itemize}
}
\slide{Teilprozesse des Gedächtnisses (Computer-Metapher)}{
    \begin{itemize}
        \item \b{Enkodierung:} Einschreiben von Informationen in das Gedächtnis (Analogie: Daten über die Tastatur eingeben).
        \item \b{Speicherung:} Dauerhaftes Behalten der Informationen (Analogie: Sichern der Daten auf der Festplatte/Server).
        \item \b{Abruf (Retrieval):} Wiedergeben von Informationen aus dem Speicher (Analogie: Öffnen einer Datei auf dem Bildschirm).
    \end{itemize}
}
\slide{Sensorischer Speicher \& Partial Report Procedure}{
    \begin{itemize}
        \item Eigenschaften: Speichert ein exaktes, unverzerrtes Abbild des Stimulus (ikonisch/echoisch) für sehr kurze Zeit (ca. 250ms).
        \item \b{Partial Report Procedure} (Sperling):
        \begin{itemize}
            \item Ablauf: Matrix aus Buchstaben wird kurz gezeigt, gefolgt von einem Ton (hoch/mittel/tief), der anzeigt, welche Zeile berichtet werden soll.
            \item Befund: Erfolgt der Ton unmittelbar nach der Präsentation, können fast alle Buchstaben der geforderten Zeile berichtet werden (Indiz für großes, aber flüchtiges Speicherabbild).
            \item Vergessen: Verzögert man den Ton (über 250ms), sinkt die Leistung drastisch ab, da der sensorische Speicher verblasst ist.
        \end{itemize}
    \end{itemize}
}
\slide{Kurzzeitgedächtnis-Modell (Atkinson \& Shiffrin)}{
    \begin{itemize}
        \item Funktion: Dient als \it{Working Memory} zur kurzfristigen Speicherung und Koordination.
        \item Enkodierung: Nimmt Informationen aus dem sensorischen Speicher auf.
        \item Begrenzung: Hat eine begrenzte Kapazität (Millersche Zahl) und Haltedauer (wenige Sekunden ohne Rehearsal).
        \item Interaktion mit LZG: Koordiniert den Transfer ins Langzeitgedächtnis (durch Rehearsal) und den Abruf von dort.
        \item Vergessen: Informationen gehen durch Zerfall (Decay) oder Interferenz verloren, wenn Speicherplätze für Neues benötigt werden.
    \end{itemize}
}
\slide{Primacy- und Recency-Effekt (Listenlängenexperiment)}{
    \begin{itemize}
        \item Experiment: Freie Wiedergabe einer gelernten Wortliste.
        \item \b{Primacy-Effekt:} Wörter am Anfang der Liste werden besser erinnert, da sie durch Rehearsal bereits öfter wiederholt und ins Langzeitgedächtnis übertragen wurden.
        \item \b{Recency-Effekt:} Wörter am Ende der Liste werden besser erinnert, da sie sich noch im Kurzzeitgedächtnis befinden und direkt abgerufen werden können.
        \item Mittelteil: Wörter in der Mitte werden am schlechtesten erinnert (weder LZG noch KZG optimal genutzt).
    \end{itemize}
}
\slide{Brown-Peterson-Paradigma}{
    \begin{itemize}
        \item Ziel: Untersuchung der Behaltensdauer im KZG ohne Wiederholung (Rehearsal).
        \item Ablauf:
        \begin{enumerate}
            \item Einprägen von drei Buchstaben (Trigramm)
            \item Rückwärtszählen in 3er-Schritten für eine variierende Zeitspanne (Retentionsintervall)
            \item Wiedergabe der Buchstaben
        \end{enumerate}
        \item Ergebnis: Die Erinnerungsleistung sinkt nach wenigen Sekunden rapide ab (nach 18 Sek. fast null).
        \item Artikulatorische Suppression: Das Rückwärtszählen belegt das verbale System ("innere Stimme") und unterdrückt so den Rehearsal-Prozess. Dies beweist, dass Informationen im KZG ohne aktive Wiederholung sehr schnell zerfallen.
    \end{itemize}
}